\chapter{Conclusion}
\label{ch:conclusion}
\section{Summary}
The identification of novel drug-target interactions remains a critical bottleneck in early-stage drug discovery, where computational approaches offer significant potential to reduce both cost and time.
This project developed a complete pipeline for DTI prediction.
It starts by preprocessing raw datasets and engineering features with industry standard encodings, and leads to an iterative development of a deep neural network.
The pipeline was made accessible with a React-based interface through FastAPI backend, providing a virtual screening interface designed for use in the early stages of drug discovery.

\section{Objectives}
\label{sec:objectives}
This section evaluates the project against the three gaps identified in the gap analysis.
\subsection{Competing with State-of-the-Art Models}
\label{subsec:competing-with-state-of-the-art-models}
The final model competes directly with current State-Of-The-Art prediction methods in DTI.
When evaluated with an increased sample size, the model achieved an AUROC of approximately 0.97, exceeding both MINDG (~0.95) and DeepCDA (~0.89) across all metrics we report.
It should be noted that the evaluation methodology differs, meaning the comparison is approximate rather than direct.
Nonetheless, the results demonstrate that the architecture is capable of competitive performance.
\subsection{Addressing Class Imbalance}
\label{subsec:addressing-class-imbalance}
The gap analysis identified class imbalance as a key challenge, with positive interactions being harder to identify than negatives.
We used the weighted loss function BCEWithLogitsLoss in PyTorch to address this, which increased the penalty for misclassifying positive interactions during training.
This allowed the model to learn hard-to-identify binding patterns rather than defaulting to predicting the majority class.
\subsection{Closing the Usability Gap}
\label{subsec:closing-the-usability-gap}
The React-based frontend provides an intuitive interface where users can select drugs and protein targets in plain English rather than requiring raw SMILES strings or protein sequences.
This directly addresses the gap identified in the literature, where most high-performing DTI models lack accessible interfaces for practical use, thus decreasing the barriers for use of the model.

\section{Contributions}
This project makes three key contributions.
First, it proves that a complex deep learning model for DTI prediction can be made accessible through an intuitive web application, closing the gap between machine learning pipelines and practical usability for researchers in early-stage drug discovery.
Second, an iterative development methodology which shows incremental improvements and provides transparency on how architectural, training, and framework decisions impact model performance.
Third, Using a recent version of BindingDB, November 2025, demonstrates its ability to work on up-to-date data and give accurate results.

\section{Reflection}
\label{sec:reflection}
This project introduced new concepts across the full pipeline of developing and deploying a deep learning system, from raw data processing to a user-facing application.
Iterating from NumPy, building the backpropogation function from scratch, to PyTorch provided insight into understanding the mechanics that frameworks abstract away.
This made the subsequent migration to PyTorch far more informed, as decisions around optimiser selection, weight initialisation, and regularisation could be made with an understanding of their underlying effects.
Developing an intuitive frontend and connecting it to the model using a RESTful API was challenging, requiring significant work in API design and state management that extended beyond the machine learning focus of the project.
Adopting cross-validation rather than a fixed data split could have further improved the performance and strengthened the evaluation when comparing it against state-of-the-art models.
Overall, the project developed skills across machine learning, backend engineering, and frontend development that extend well beyond the immediate scope of DTI prediction.

